\chapter{科研辅助分析平台的集成设计与应用验证}
\enchapter{Integrated Design and Application Validation of the Scientific Research Assistant Analysis Platform}
\label{chap:platform_design}

\section{引言}
\ensection{Introduction}

前文围绕碳纳米管阵列（CNT Array）研究中的关键分析任务，分别开展了知识建模与检索增强、SEM 图像形貌表征以及工艺参数与形貌关系建模等方面的研究，并取得了相应成果。其中，第三章面向科研知识利用需求，构建了适用于 CNT 阵列领域的知识组织与检索增强方法，为相关问题的知识查询、关联分析与证据获取提供了支撑；第四章则从实验表征与工艺分析角度出发，分别实现了 SEM 图像形貌特征的自动提取与量化表征，以及工艺参数驱动的关键形貌指标预测与关系解释。上述研究分别在特定任务中验证了方法的有效性，为 CNT 阵列科研分析提供了方法基础。

然而，从实际科研工作流程来看，单一方法的独立应用仍难以满足研究人员在样品分析、工艺优化与结果解释中的综合需求。科研活动往往需要将多源数据管理、知识检索分析、图像自动表征以及预测推断等能力协同起来，形成面向具体研究对象和实际问题的统一分析工具。因此，有必要在前述研究成果的基础上，进一步开展系统集成设计，将相关方法部署到统一的平台环境中，构建面向 CNT 阵列研究的科研辅助分析平台，以提升方法的实用性、协同性与可用性。

基于此，本章以前述研究成果的工程化落地与综合应用为目标，重点研究科研辅助分析平台的总体架构、模块实现与部署方式。围绕 CNT 阵列研究中的知识辅助分析、SEM 图像自动表征以及工艺参数驱动的形貌预测与解释等典型任务，构建集数据管理、算法调用与结果展示于一体的平台体系，并建立第三章、第四章相关方法与系统功能模块之间的映射关系。在此基础上，通过典型应用场景对平台的实际运行效果进行验证，分析其在科研流程支撑、结果输出与应用可行性等方面的表现。

本章的研究内容不仅是对前文方法成果的系统整合，也是对其应用价值的进一步验证。通过平台化实现，可将知识分析、图像表征与预测解释等能力纳入统一科研辅助流程，从而为 CNT 阵列样品分析、工艺优化与实验结果理解提供更为完整的技术支撑。

\section{系统总体架构与部署设计}
\ensection{Overall Architecture and Deployment Design}

\subsection{系统建设目标}
\ensubsection{System Construction Objectives}

面向 CNT 阵列研究的科研辅助分析平台，并非追求通用办公或通用数据分析场景下的软件功能完备性，而是围绕领域科研流程中的实际问题构建专用系统。结合前文研究内容以及实验分析需求，平台建设目标主要集中在知识辅助分析、SEM 图像自动表征以及工艺参数驱动的形貌预测与解释三类典型任务上。其核心任务不是替代研究人员完成科研判断，而是在数据整合、算法调用、结果组织和证据支撑等环节提供系统化辅助，以提高科研分析效率并增强结果的可追溯性。

首先，在知识辅助分析方面，平台需要能够围绕科研问题组织领域知识，为研究人员提供与工艺调控、形貌变化和机理解释相关的结构化证据。与面向开放域问答的检索系统不同，本研究中的知识分析更强调领域约束、样品上下文和关系链条的完整性。因此，平台需要支持围绕科研问题开展查询解析、知识召回、关系扩展和证据组织，使研究人员能够从平台输出中直接获取与当前问题相关的分析依据。

其次，在 SEM 图像自动表征方面，平台应能够对实验获取的图像数据进行统一接入和自动处理，将原始图像转化为可用于分析的关键形貌指标。对于 CNT 阵列研究而言，SEM 图像是连接实验观察与后续工艺分析的重要媒介。若形貌特征仍主要依赖人工判读，不仅效率较低，也不利于样品间比较和结构化归档。因此，平台需要具备图像上传、预处理、模型推理以及形貌指标结构化输出能力，使图像分析能够稳定嵌入科研流程。

再次，在工艺参数驱动的预测与解释方面，平台应能够根据输入工艺条件输出关键形貌指标的预测结果，并进一步结合知识检索结果补充解释信息。单纯给出预测值虽然能够支持工艺比较，但难以满足科研场景中对可解释性和证据支撑的要求。因此，平台建设目标还包括形成“预测结果 + 知识证据”的联合分析模式，使模型输出能够与领域知识共同支撑实验设计与结果判断。

综合而言，平台建设的总体目标可以概括为三个方面：一是实现前文提出方法的系统化部署，使知识检索、图像表征和关系建模能力能够在统一环境中协同运行；二是围绕样品分析、工艺优化与结果解释建立贯通的数据与业务链路，提升科研分析流程的连贯性；三是通过集成化界面和结构化结果组织，提高前述方法在实际科研场景中的可用性、可追溯性与应用价值。

\subsection{系统总体架构}
\ensubsection{Overall System Architecture}

为满足上述建设目标，平台整体采用分层架构设计，将系统划分为前端展示层、后端业务层、算法服务层和数据存储层四个层次。该分层方式一方面有利于将界面交互、业务逻辑与算法计算解耦，另一方面也便于将第三章和第四章提出的方法以模块化方式部署到统一平台中，从而形成较清晰的系统结构。

前端展示层主要负责用户交互、任务提交和结果展示，是研究人员直接使用平台的入口。该层围绕科研分析流程提供样品检索、问题输入、图像上传、参数录入以及结果查看等交互界面，并对知识证据、形貌指标和预测结果进行可视化展示。由于平台面向科研分析任务而非普通信息检索任务，前端不仅承担输入输出功能，还需要支持样品上下文切换、结果追溯和不同模块之间的联动展示，以保证研究人员能够在统一界面中查看同一样品相关的多类分析结果。

后端业务层主要负责请求调度、任务管理、结果封装以及模块间协同，是连接前端界面与底层算法服务的核心枢纽。当用户发起知识查询、图像表征或参数预测任务时，后端根据任务类型组织对应的数据读取、参数校验和算法调用流程，并将多模块输出结果进一步整理为统一的数据结构返回给前端。对于需要跨模块协同的任务，例如基于预测结果补充知识解释或围绕样品编号汇总图像与工艺信息的场景，后端业务层还承担结果拼接与上下文组织功能。

算法服务层是平台实现科研分析能力的关键部分，主要部署前文提出的知识检索增强方法、SEM 图像自动表征方法以及工艺参数驱动的预测模型。具体而言，第三章提出的知识建模与检索增强方法在系统中对应知识检索与分析模块，负责围绕科研问题完成查询解析、混合召回、关系约束扩展和证据链输出；第四章提出的 SEM 图像自动分析方法在系统中对应图像表征模块，负责图像预处理、前景分割和关键形貌指标提取；第四章中工艺参数与形貌关系建模方法则在系统中对应预测与解释模块，负责根据输入工艺参数给出关键形貌指标预测结果，并与知识检索模块联动生成解释信息。通过这样的映射关系，前述方法不再作为独立实验流程存在，而是被封装为可供平台调度的算法能力。

数据存储层主要用于统一管理平台运行过程中涉及的多源数据，包括样品基础信息、工艺参数记录、SEM 图像文件、结构化知识单元、形貌特征结果、预测结果以及分析日志等内容。与一般业务系统相比，本研究中的数据组织更强调样品级绑定关系，即以样品编号为主线组织不同来源的数据对象，使知识分析、图像表征和预测结果之间能够围绕同一研究对象建立对应关系。这样既有利于模块之间共享上下文，也为结果回写、历史追溯和后续比较分析提供了基础。

总体来看，该四层架构既能够清楚划分系统各层功能边界，又能够将第三章和第四章的方法成果映射为平台中的可调用模块，从而实现科研辅助分析能力的统一集成。前端展示层负责人机交互，后端业务层负责流程组织，算法服务层负责核心分析能力实现，数据存储层则提供持续稳定的数据支撑，四者共同构成面向 CNT 阵列研究的科研辅助分析平台。

\subsection{系统部署与运行流程}
\ensubsection{System Deployment and Operation Workflow}

为了提高平台的可维护性与功能扩展能力，系统在实现上采用前后端分离并结合服务化调用的部署方式。前端界面负责接收用户输入并展示分析结果，后端系统则通过统一接口调度知识检索、图像处理和预测模型等核心能力。对于算法功能部分，不同模块以相对独立的服务形式接入后端业务层，从而避免将复杂模型推理流程直接耦合在界面层或单体应用逻辑中。该部署方式既便于后续替换或扩展算法模型，也有利于根据任务特点调整不同模块的运行环境。

在知识检索能力接入方面，系统将第三章构建的知识库、索引结构及关系约束检索流程封装为可调用服务。当前端提交科研问题后，后端首先解析查询内容与任务类型，再将相应请求发送至知识检索模块完成混合召回、关系扩展与证据链组织，最终返回结构化答案、关键证据片段以及相关来源信息。该过程使知识分析结果不仅可作为独立输出，也可为后续形貌解释与预测解释任务提供支撑。

在图像处理能力接入方面，系统将第四章构建的图像自动表征流程作为独立处理链路部署。用户上传 SEM 图像后，后端首先完成样品编号绑定、文件存储和基础校验，再调用图像表征模块执行预处理、前景分割、骨架与几何分析及关键形貌指标提取。得到的结构化结果随后写入数据库，并可由前端进一步展示为指标列表、分析报告或样品级结果摘要。

在预测模型接入方面，系统将以工艺参数为输入、以关键形貌指标为输出的预测能力封装为独立服务。当前端录入或选定工艺参数组合后，后端完成参数格式化与范围校验，并调用预测模块输出相应形貌指标预测值。在此基础上，若用户需要对预测结果进行进一步解释，系统还可将预测目标、核心参数和形貌指标作为检索上下文传递给知识检索模块，生成与预测结果相对应的知识证据。

从整体业务链路看，平台运行流程可以概括为“输入数据或问题提交、后端任务分发、算法模块处理、结果组织回写、前端统一展示”五个阶段。对于知识辅助分析任务，流程以用户问题为起点，经知识检索模块处理后输出分析结论与证据链；对于图像自动表征任务，流程以 SEM 图像上传为起点，经图像表征模块处理后输出关键形貌指标；对于预测与解释任务，流程以工艺参数输入为起点，经预测模块和知识检索模块联动处理后输出预测结果及其解释信息。通过上述部署与运行方式，平台形成了从输入到分析再到结果反馈的完整业务闭环，为后续典型应用场景验证奠定了基础。

\section{核心功能模块实现}
\ensection{Implementation of Core Functional Modules}

\subsection{多源数据管理模块}
\ensubsection{Multi-source Data Management Module}

科研辅助分析平台的稳定运行依赖于对多源数据的统一组织与管理。在 CNT 阵列研究场景中，平台需要同时处理样品基础信息、工艺参数、SEM 图像、知识单元以及分析结果等不同类型的数据对象。若这些数据彼此分离管理，不仅难以在模块之间共享上下文，也会削弱结果追溯和样品比较分析的可行性。因此，平台首先构建多源数据管理模块，对不同来源的数据进行统一接入、绑定与维护。

该模块以样品编号为核心组织主线，将同一样品对应的工艺记录、图像数据、形貌特征和预测结果关联到统一的数据对象中。样品编号不仅用于存储层中的主键关联，也贯穿平台前端检索、后端调度和分析结果回写全过程。通过以样品为中心组织数据，平台能够在知识检索、图像分析和预测任务之间建立共享语境，避免不同模块在处理同一研究对象时出现信息割裂。

在数据类型上，多源数据管理模块主要包括四类内容。第一类是样品基础数据与工艺参数数据，用于记录样品来源、实验批次和制备条件；第二类是 SEM 图像及其元数据，用于保存图像文件路径、采集条件及与样品的绑定关系；第三类是知识单元数据，用于存储第三章构建的结构化知识条目及其来源信息；第四类是分析结果数据，包括图像表征得到的关键形貌指标、预测模块输出的预测结果以及知识检索结果的证据链记录等。上述数据通过统一的数据模式和字段约束进行管理，从而为上层功能模块提供规范输入。

除基础存储外，该模块还承担数据预处理和结果回写功能。对于新接入的图像文件或实验记录，系统会在写入前执行字段校验、编号匹配和必要的格式转换，以保证后续算法模块能够稳定调用。对于平台运行过程中产生的形貌指标、预测结果和知识证据，系统则会在任务完成后将结果回写至相应数据表，并保留时间戳、任务类型和来源信息，以支持结果追溯、历史比较和样品级分析汇总。

因此，多源数据管理模块在整个系统中并非简单的数据存储单元，而是支撑知识检索、图像分析、预测计算及结果组织的基础模块。其核心价值在于通过样品级数据绑定与结构化管理，为科研辅助分析平台提供统一、可追溯且可联动的数据基础。

\subsection{知识检索与分析模块}
\ensubsection{Knowledge Retrieval and Analysis Module}

知识检索与分析模块主要承接第三章提出的知识建模与检索增强方法，是平台中面向科研问题分析的重要组成部分。该模块的目标不是提供开放式通用问答，而是围绕 CNT 阵列研究中的工艺分析、形貌解释和预测解释等任务，输出与问题相关且具有来源支撑的结构化证据结果。为此，模块在实现过程中保留了第三章提出的知识单元组织方式与关系约束检索机制，并将其封装为适用于系统调用的业务流程。

当用户在平台中输入科研问题后，系统首先对问题进行解析，识别其中涉及的核心对象、工艺变量、形貌指标及任务类型。例如，当问题聚焦于某一工艺参数对形貌的影响时，模块将其归类为工艺分析任务；当问题围绕已观测到的形貌现象追溯成因时，则归类为形貌解释任务。该步骤的主要作用是为后续召回范围控制和关系扩展策略选择提供依据。

在查询解析完成后，模块执行混合召回过程。该过程综合利用关键词匹配、结构化字段筛选和语义相似度检索，从知识库中召回与当前问题相关的候选知识单元。与仅依赖文本相似度的方式相比，混合召回能够更好地兼顾领域术语匹配和语义关联，从而提高科研问题场景下的初始候选质量。对于与样品编号、特定工艺条件或目标形貌指标明确相关的查询，系统还会优先利用结构化字段约束缩小候选范围。

在候选知识单元获取后，模块进一步依据第三章构建的关系类型和约束规则进行关系扩展与证据组织。该过程围绕问题中的核心变量、目标对象和限制条件，对候选知识单元进行链式连接与筛选，形成可用于说明“工艺参数变化 - 机理作用 - 形貌结果”或“预测目标 - 相关工艺 - 支撑证据”等关系的证据链。相较于简单返回若干相似文本片段，该过程更强调证据之间的逻辑衔接与上下文一致性，因而更适合科研分析任务。

最终，知识检索与分析模块输出的结果通常包括问题对应的分析结论、关键证据片段、来源信息以及必要的关系链说明。对于平台中的独立知识辅助分析任务，这些内容可以直接展示为问题答案与证据支持；对于图像分析解释或预测结果解释任务，这些输出则进一步作为附加说明与对应结果联动展示。由此可见，该模块不仅实现了第三章方法的系统落地，也在平台中承担了科研知识支撑和解释增强的重要角色。

\subsection{图像表征模块}
\ensubsection{Image Characterization Module}

图像表征模块主要对应第四章提出的 SEM 图像自动分析方法，负责将实验获取的原始 SEM 图像转化为可直接用于比较分析和建模的结构化形貌指标。对于 CNT 阵列研究而言，SEM 图像不仅承载样品形貌信息，也是连接工艺条件与结构结果的重要中间表征。因此，在平台中实现稳定的图像自动表征流程，是支撑样品级分析和后续预测建模的关键环节。

在具体实现上，图像表征模块的处理流程包括图像上传、预处理与分割、特征提取以及结构化结果输出四个阶段。首先，用户通过平台上传 SEM 图像，系统对图像文件进行基础校验，并完成样品编号绑定、文件存储与元数据登记。随后，模块执行图像预处理，包括尺寸统一、主体区域提取及必要的灰度增强，以提高后续模型推理和几何分析的稳定性。

在预处理完成后，模块调用第四章中构建的前景分割模型，对图像中的 CNT 前景区域进行提取。与论文前文中的离线实验流程相比，平台化实现更强调模块调用的稳定性与结果可回写性，因此在系统中采用封装后的标准化推理接口输出二值前景图和中间处理结果。得到稳定前景后，模块进一步执行骨架分析与几何统计，提取覆盖密度、取向度、表观直径以及平均曲率、波曲度、迂曲度等关键形貌指标，并在样品级层面完成结果聚合。

模块输出的结果以结构化方式写入数据存储层，并通过前端界面展示为可供浏览和比较的形貌指标列表。必要时，平台还可以同步显示原始图像、分割结果和核心特征值，以帮助研究人员快速核验自动分析结果。通过这种方式，图像表征模块不仅实现了关键形貌特征的自动提取与量化表达，也将 SEM 图像分析从独立的离线处理流程转化为平台内可直接调用的科研辅助功能。

\subsection{预测与解释模块}
\ensubsection{Prediction and Interpretation Module}

预测与解释模块主要用于将工艺参数与关键形貌指标之间的关系建模结果转化为平台中的可用功能。该模块以用户输入的工艺参数为基础，输出目标样品在给定条件下的关键形貌指标预测值，并进一步结合知识检索模块提供相应解释信息。其目标不仅在于给出数值预测结果，更在于通过知识证据补充预测背景，从而增强结果的可解释性和科研参考价值。

在预测部分，模块首先接收用户输入或样品选择带入的工艺参数，包括与催化层制备、退火过程和生长过程相关的核心变量。系统对输入参数进行格式化和有效性检查后，将其传递给训练完成的预测模型，输出对应的关键形貌指标预测结果。结合第四章研究内容，平台中的预测对象重点包括取向度、平均曲率、波曲度和迂曲度等能反映阵列结构状态的指标。预测结果可直接用于不同工艺组合之间的趋势比较和实验方案辅助分析。

在解释部分，模块并不将预测结果视为孤立数值，而是进一步调用知识检索与分析模块，围绕输入工艺变量、目标形貌指标及其变化方向组织相关证据。具体而言，系统将预测任务中涉及的核心参数和目标指标转换为知识检索上下文，从知识库中召回与之相关的工艺--形貌关系、机理说明和已有实验观察结果，并进一步整理为与当前预测结果相对应的解释信息。这样，平台最终输出的不仅是预测值本身，还包括与该结果相关的知识依据和解释线索。

因此，预测与解释模块形成了“预测结果 + 知识证据”的联合输出模式。一方面，预测模型为工艺优化和实验设计提供定量参考；另一方面，知识检索结果为预测值的合理性提供定性支撑，使研究人员能够在查看结果时同时获得与之对应的工艺背景和机理说明。对于科研辅助分析平台而言，这种联合输出方式比单一预测结果更符合科研分析中对解释性和可追溯性的要求。

\section{系统实现与应用效果}
\ensection{System Implementation and Application Performance}

\subsection{系统实现与运行环境}
\ensubsection{System Implementation and Runtime Environment}

平台实现采用前后端协同开发方式，其中前端主要负责任务交互与结果展示，后端主要负责业务逻辑组织与算法模块调度，数据层用于保存实验数据与分析结果，算法层用于承载知识检索、图像分析与预测模型等核心能力。通过这种分工，系统能够在统一界面中集成不同类型的科研辅助能力，并支持围绕样品分析流程开展数据管理与结果联动。

在运行环境方面，平台具备完成知识检索、图像处理和预测计算所需的基本条件。知识检索模块依赖结构化知识库、索引服务与后端接口协同运行；图像表征模块运行于支持模型推理和图像处理的环境中；预测模块则依赖已训练完成的形貌预测模型及其推理接口。系统运行过程中，前端界面通过网络请求与后端交互，后端根据任务类型调用对应算法服务，并将结果统一写入数据库后返回前端展示。由此，平台能够形成完整的模块部署与功能调用条件，满足第五章所讨论的系统集成与应用验证需求。

需要说明的是，本章重点在于验证前文方法的系统集成能力与应用可行性，因此系统实现部分更强调平台能够稳定支撑相关功能模块运行，而非对通用软件工程细节作展开讨论。就当前实现情况来看，平台已具备样品数据接入、知识查询、图像自动表征和参数预测分析等功能的部署基础，能够支持后续典型应用场景的运行展示。

\subsection{典型应用场景展示}
\ensubsection{Typical Application Scenarios}

为验证平台在 CNT 阵列研究中的实际应用效果，本文结合平台面向科研流程的功能定位，选取知识辅助分析、SEM 图像自动表征以及工艺参数驱动预测与解释三个典型应用场景进行展示。三个场景分别对应第三章与第四章提出方法在平台中的落地应用，也能够较完整地覆盖科研分析中“问题分析 - 实验表征 - 结果预测与解释”的主要流程。

在知识辅助分析场景中，研究人员以工艺调控或形貌变化相关问题作为输入，平台通过知识检索与分析模块对问题进行解析，并从知识库中返回相关证据链和解释结果。该场景的输入主要为科研问题文本，系统处理过程包括问题解析、候选召回、关系约束扩展与证据组织，输出则表现为与问题相关的分析结论、关键证据以及来源信息。该场景表明，平台能够将第三章中的知识检索增强方法转化为面向科研问题的直接辅助工具，为工艺分析和形貌成因判断提供知识支撑。

在 SEM 图像自动表征场景中，研究人员上传具体样品的 SEM 图像，平台自动完成图像预处理、前景分割和形貌特征提取，最终输出对应样品的关键形貌指标。该场景的输入为原始图像数据，系统处理过程包括图像文件接入、自动分析流程调用以及特征聚合，输出则表现为覆盖密度、取向度、表观直径和弯曲相关特征等结构化结果。该场景说明，平台能够将第四章中的自动化形貌表征方法嵌入实际分析流程，为实验样品的快速量化比较提供支撑。

在工艺参数驱动的预测与解释场景中，研究人员输入目标工艺参数组合，平台首先调用预测模型给出关键形貌指标的预测值，随后再联动知识检索模块返回与预测结果相关的解释证据。该场景的输入为工艺参数或候选工艺组合，系统处理过程包括参数校验、模型推理和知识证据补充，输出则表现为预测指标结果及其对应解释信息。与仅输出预测值的方式相比，该场景体现了平台在结果解释与科研参考方面的进一步能力，使研究人员能够同时获得定量预测和定性支撑。

综合这三个典型场景可以看出，平台并非简单地将多个独立算法并列展示，而是围绕 CNT 阵列研究中的实际分析需求，将知识、图像和预测能力组织为可服务科研流程的统一工具。其应用意义在于，一方面能够减少研究人员在数据整理、方法切换和结果归纳上的重复工作，另一方面也增强了样品分析、工艺比较和结果解释过程中的系统性与连贯性。

\subsection{系统运行效果分析}
\ensubsection{System Performance Analysis}

从系统集成角度看，平台较好地实现了前文研究方法的统一部署。第三章提出的知识建模与检索增强方法在平台中转化为知识检索与分析模块，第四章提出的图像自动表征方法和工艺参数与形貌关系建模方法则分别转化为图像表征模块和预测与解释模块。通过统一的数据组织和后端调度机制，这些方法不再作为彼此分离的实验流程存在，而是能够围绕同一研究对象和同一分析任务协同工作。因此，平台在方法集成程度方面达到了将前文研究成果工程化落地的基本目标。

从科研流程支撑能力看，平台已经具备围绕样品分析、工艺优化和结果解释开展完整辅助分析的条件。多源数据管理模块提供样品级数据组织基础，知识检索模块支持问题分析与解释生成，图像表征模块支持样品形貌自动量化，预测与解释模块则进一步实现工艺参数驱动的预测分析与知识补充。各模块在功能上相互衔接，使平台能够覆盖科研过程中较为典型的几个关键环节。尤其是在需要对实验观察结果进行快速解释或对候选工艺组合进行辅助判断时，平台的系统化组织优势较为明显。

从结果展示效果看，平台输出结果具有一定的结构化与可读性特征。知识分析结果以证据链和来源信息形式呈现，有助于研究人员理解分析依据；图像表征结果以关键形貌指标形式输出，便于不同样品间比较；预测模块在输出预测值的同时补充知识证据，也增强了结果解释性。相较于仅返回单一数值或分散文本片段的形式，平台更有利于在科研语境中组织结果和支持后续讨论。

从运行效率与使用便利性看，平台通过统一界面完成问题输入、图像上传、参数录入和结果查看，降低了研究人员在不同工具之间切换的成本。与此同时，模块化部署方式也使不同算法功能能够在后端按需调用，避免科研分析流程过度依赖人工文件整理和离线脚本切换。当然，平台当前仍以支撑论文研究验证为主要目标，在大规模并发访问、复杂权限管理和更广泛实验体系适配等方面仍有进一步完善空间，但这并不影响其作为专用科研辅助分析平台在当前应用场景中的可行性与有效性。

\section{小结}
\ensection{Summary}

本章围绕科研辅助分析平台的集成设计与应用验证，构建了面向 CNT 阵列研究的系统总体架构与部署方案，并设计实现了多源数据管理、知识检索与分析、图像表征以及预测与解释等核心功能模块。在此基础上，通过典型应用场景对平台运行效果进行了验证。研究结果表明，前文提出的知识增强、形貌表征与关系建模等方法不仅能够在单项任务中发挥作用，而且具备进一步集成为科研辅助分析平台的系统落地能力，可为 CNT 阵列样品分析、工艺优化与结果解释提供一体化支撑。
